\documentclass [varwidth]{standalone}
%scale[1]
%lmargin[30]
%vmargins[30]
\usepackage {texmaths-main}
\begin {document}
\setCompleteVersion {0}\def \disquesize {0.8}\noindent \begin {tblr}{colspec={Q[wd=63pt]Q[wd=43pt]Q[wd=43pt]Q[wd=43pt]Q[wd=88pt]},hlines,vlines} Pour représenter des {\bfseries demis} de ce disque, on le partage en ....... parties égales.&\disque \disquesize \par \kern 10pt$\dfrac 12=$&\disque \disquesize \par \kern 10pt $\dfrac 22=$&\SetCell [c=2]{l}\disque \disquesize \kern 10pt\disque \disquesize \par \kern 10pt$\dfrac 32=$\par \kern 5pt&\\ Pour représenter des {\bfseries tiers} de ce disque, on le partage en ....... parties égales.&\disque \disquesize \par \kern 10pt$\dfrac 13=$&\disque \disquesize \par \kern 10pt $\dfrac 23=$&\disque \disquesize \par \kern 10pt$\dfrac 33=$&\disque \disquesize \kern 10pt\disque \disquesize \par \kern 10pt$\dfrac 43=$\par \kern 5pt\\ Pour représenter des {\bfseries quarts} de ce disque, on le partage en ....... parties égales.&\disque \disquesize \par \kern 10pt$\dfrac 14=$\par \kern 10pt&\disque \disquesize \par \kern 10pt $\dfrac 24=$&\disque \disquesize \par \kern 10pt$\dfrac 44=$&\disque {0.6}\kern 5pt\disque {0.6}\kern 5pt\disque {0.6}\par \kern 10pt$\dfrac 94=$\par \kern 5pt\\ \end {tblr}\par 
\end {document}
